\chapter{Review of Numerical Simulations}
\label{ch_2}
Modern cosmology aims to combine the observations of the Universe with the mathematical description provided by the Standard Model. However, all assumptions and theories must be tested. Many predictions made by the Standard Model are not always directly observable, such as predictions about Inflation, but there must be a way to verify whether the ideas of the Standard Model are self-consistent and match up with observation. This is difficult and essentially requires a computational effort. In the context of this thesis it is appropriate to look at the computational efforts that have been developed to study LSS.

Galaxy clustering is a non-linear process where the mean density of a cluster or galaxy is thousands or millions times denser than the mean background density \citep{bib_pea}. As the evolution of such a system is inherently non-linear and in general lacks an analytic solution then one must resort to `brute force' methods of analysis via a computer. The idea here is to use a computer simulation to study the evolution of one possible realization that should be statistically equivalent to the real Universe. Such an idea is similar to exhaustive proofs in mathematics where an elegant proof is non-trivial. One such method of simulation is called an $N$-body simulation: it follows the motion of many particles that are guided by the physical laws encoded in the simulation. Observations and physical intuition provide constraints on the parameter space that such a simulation will explore. The end result should be a statistical equivalence between the simulated Universe and the observed one.

Bertschinger reviews the main concepts and history of cosmological simulations in his 1998 paper: ``Simulations of structure formation in the Universe'' \citep{bib_bert}. This paper covers the underlying theory and offers insight into each of the various methods of performing LSS simulations. In this paper he says that the first gravitational simulation was performed on an analogue optical computer by Holmberg in 1941, where lightbulbs and photo detectors were used to replicate the inverse square law of gravity. The first digital gravitational simulations were by von Hoerner in 1960 and 1963, and Aarseth in 1963. Bertschinger attributes the first cosmological simulation to Press and Schechter in 1974; what they did was to investigate the mass distribution of bound clumps formed by hierarchical clustering. They constructed a model for predicting the number of objects of a certain mass within a given volume of the Universe, and has come to be known as the ``Press-Schechter formalism" \citep{bib_pres} (also known as the cloud-in-cloud problem).

In the thirteen years since Bertschinger's review paper further progress has been made and one of the first billion-body ($N = 1024^3$) gravitational simulations (Millennium Simulation \citep{bib_spr}) was carried out in the year 2005. The code used (GADGET-2) in this simulation was a mix of the direct summation (particle-particle, P-P) and particle mesh methods, sometimes denoted as a ${\rm P^3 M}$ code, it also uses the Tree and SPH methods (omitted in this thesis for brevity). Both P-P and P-M methods are outlined below in sections \ref{sec.direct} and \ref{sec.mesh}.

By today's standards, running an LSS code with $N = 512^3$ resolution is considered coarse \citep{bib_hmpc}. The latest simulations have a resolution of about one hundred billion ($10^{11}$) particles, the RAMSES code by \citet{bib_tess} claimed to simulate LSS using 70 billion particles. The number of particles, denoted by $N$, in an $N$-body code (or the number of grid-points in fluid/mesh codes) dictates the resolution of the simulation; more particles allows for higher resolution.

As mentioned in the previous chapter there are alternatives to the $N$-body method. Large Scale Structure can also be simulated as a hydrodynamical pressureless fluid whereby each particle, or mass-containing fluid element, may represent a galaxy.

For computational purposes, the most natural approach is to treat each galaxy as a particle in the simulation and to directly calculate the gravitational force between each particle. The computational particles, here, only account for the gravitational force all other physics is essentially smoothed out. It is more accurate to say that the computational particle represents collection of real particles that have the same mass as a galaxy. \citet{bib_he} say that the computational particle is a superparticle: it represents many real particles, where the underlying physics is either unknown or otherwise not included in the simulation.

Calculating the direct gravitational force between all such superparticles in a simulation is perhaps the slowest method and scales as $N \times N$. This is not the most efficient way but a variant of this algorithm is still used. Here is a mostly complete list of the standard computational methods:
\begin{itemize}
\item Direct Force calculation of Particle-Particle interactions
\item Particle Mesh and ${\rm P^3 M}$
\item Tree Codes
\item Fluid approaches
\item Phase space methods
\end{itemize}

\noindent In the following sections I provide a review of the first two techniques on this list. The last two techniques are mentioned in the last section \ref{lsswm_con} of this chapter but further discussion of tree codes is omitted; however, details on tree codes can be found in Bertschinger (\citeyear{bib_bert}).

\section{$N$-Body simulations}
\label{sec.dir_nb}
%slow - fully non-linear
%softening / truncation
The $N$-Body problem considers $N$ point masses that under go classical gravitational attraction. Such a system is deceptively simple but ultimately has many caveats that need to be considered. $N$-body systems have a tendency to be chaotic, only for simulations of two-bodies can there be a trivial general solution. In the case of celestial mechanics, stability and chaos are central themes \citep{bib_aar}. In a simulation of the solar system one particle might represent each celestial body (for example, the Sun and the planets). However, in a cosmological simulation it is less obvious what we mean by a ``particle''. \citep{bib_he, bib_hmpc}

In simulations of LSS, the point particles may represent real masses of whole galaxy clusters down to sub-galaxy sized objects. For a given box size, a high mass resolution in the simulation comes from having a large number of particles. In a simulation of a large box size but low mass resolution then one particle may represent a galaxy cluster in terms of mass. Likewise, having a large number of particles in a large box will allow for sub-galaxy masses for each particle as seen in the Millennium Simulation \citep{bib_spr}. This begs the question of ``what is a particle in a simulation?'' Clearly, one particle in the simulation represents many real particles. For this reason particles in simulations are commonly referred to as superparticles \citep{bib_he}. The more particles there are, the higher the resolution and the better the approximation gives to the underlying continuous density field (although not necessarily a better approximation to the physical processes involved). Consequently, the minimum useful mass resolution is $~10^9$ (this dimensionless ratio is, the mass range from sub-galactic scales to super-clusters) \citep{bib_hmpc}.

Couchman arrives at this resolution by using the following considerations: to model a fair sample of the Universe requires a box size of at least $(100 {\rm Mpc})^3$, this volume will contain about $10^4$ bright galaxies. Allowing for 100 particles per galaxy halo (in terms of mass) suggests that a simulation should have at least $10^8$ resolution elements or particles \citep{bib_hmpc3}.

In Newtonian physics the gravitational force between any two objects is a non-linear function of position, although the force can at least be approximated as linear when the system has barely evolved. Literature suggests that finding a solution to a general `many-body' problem is impossible, at least to the extent of only considering first integrals. A general solution for $N=3$ (Where $N$ is the number of bodies) was found by Karl Sundman in 1912 through a process known as regularization (which assumes zero angular momentum --- a fair assumption for gravity which is a radial force). The process of regularization is a way of dealing with collisions of bodies (singularities in the equations), Sundman avoided three body collisions by choice of initial conditions \citep{bib_sund}. However, it has also been shown that regularization cannot be found analytically for collisions of more than two bodies. Singularities (for example, collisions) become more complicated for $N>2$ and were omitted by Wang in his generalization of Sundman's result to $N>3$ \citep{bib_wang}. Given the difficulty of finding a general analytic solution to $N$-body problems one must perform simulations that solve the equations of motion via numerical integration.

The process of regularization will arise again in the following section on Direct Force Summation, as it deals with the singularities that can arise from particles coming too close together.

\subsection{Direct Force Summation}
\label{sec.direct}
The simplest method of performing an $N$-body simulation is to perform the direct calculation of force on each particle from every other particle. The force is calculated at every time-step and is put into the equations of motion in order to determine the new velocity and position of all the particles. Typically, this is done by using the simple Newtonian gravitational force between two particles:

\begin{equation}
\label{eqn_fgr}
\underline{F}_{\ij} = -G \sum_{j = 1,\ne i}^N m_i m_j \frac{\underline{r}_i -\underline{r}_j}{|\underline{r}_i - \underline{r}_j|^3}
\end{equation}

\noindent here $\underline{r}_i$ is the position vector of the $i$th particle. Besides cosmological simulations, which I will frequently return to in this chapter, another example of a gravitational $N$-body problem is that when the ``role of the particle model appears in the simulation of systems in which a star may be considered as a mass point with no other properties than gravitational attraction and mass" (Chapter 11, \citet{bib_he}).

One of the first examples of an $N$-body gravitational code was that of Aarseth (\citeyear{bib_aar2}), he applied his code to a cluster of galaxies with the number of particles in the range $N = 25 - 100$. Aarseth explains that ``galaxies are extended bodies but can generally be regarded as mass points. When two galaxies are involved in a collision, however, this is no longer a good approximation because the forces are not convergent."

\paragraph{Softening parameter.}
From equation \ref{eqn_fgr} we can see that as the separation between two particles decreases then the denominator of equation \ref{eqn_fgr} will tend to zero, hence the force (or, acceleration) would tend to infinity. As we approach this situation then the system will become numerically unstable. Aarseth introduced a small constant $\epsilon$ which is to be associated with the effective size of a galaxy.

\citet{bib_he} note that this $\epsilon$ parameter was soon introduced into simulations of star systems, by necessity to keep the simulation numerically stable, even though a point-mass simulation ($\epsilon = 0$) is a better physical model. These authors point out that it is not good computational practise to allow any variables to reach exceptionally large values as time integration becomes inaccurate and arithmetic overflow may occur.

To avoid these problems Aarseth introduced a regulator, that is a minimum distance ($\epsilon$) for impact $\epsilon$.
This is known as the process of regularization. In $N$-body literature this parameter is called a softening parameter. It will `smooth' out, or `soften', the computed forces. As the regulator is finite then the force is also finite \citep{bib_aar, bib_heg}.

In Chapter \ref{ch_7} I will discuss the nature of point-particles in simulations and a possible method of giving them a non-point like distribution by including higher order moments of the mass distribution. The zeroth order any mass distribution is the total mass of the particle, it has no structure or width: that is to say that it is a point. To the best of my knowledge, this is the order at which all $N$-body simulations work. The addition of a softening will be some perturbation about this zeroth order and hence provide an effective width (or size) for the particle.

For the GADGET-2 code, \citet{bib_spr} says that the single particle density function is a delta function convolved with a normalized gravitational softening kernel of comoving size $\epsilon$. He says that the particles will have a Newtonian potential of a point that is the same as Plummer sphere of size $\epsilon$.

Explicitly we can see that the addition of a softening parameter (or minimum distance) prevents the denominator of the gravitational force equation (acceleration, here) reaching zero, hence the acceleration itself is finite:

\begin{equation}
\ddot{\underline{r}}_i = -G \sum_{j = 1,\ne i}^N m_j \frac{\underline{r}_i -\underline{r}_j}{(|\underline{r}_i - \underline{r}_j|^2 + \epsilon^2)^{3/2}}
\end{equation}

\noindent $|\epsilon| << 1$ is the softening parameter. The softening parameter is non-derivable (that is, phenomenological) hence it does not appear naturally in the equations, instead it is a computational fix. This addition of a softening parameter is a phenomenological way of preventing artificial two-body relaxation. However, it is clear that without this parameter then two such particles can dominate the entire energy distribution of the system.

\citet{bib_hmpc3} suggests that the softening scale should be larger enough to avoid two-body relaxation but not too large to cause over-merger of ``fluffy'' substructures, that is why he recommends using a softening length which is an order of magnitude smaller than the interparticle separation.

A good discussion of two-body relaxation is provided in Heggie \citep{bib_heg}, here I will paraphrase the discussion:

\begin{quote}
Maxwell used the term relaxation to apply to a deformed elastic body returning to equilibrium. The idea can also be applied to the theory of gases and to stellar dynamics. However, in the latter case equilibrium is never achieved because particles can escape so it is quasi-equilibrium. In a simulation the energy of one particle is altered by its interaction with another. Normally, the interactions between two particles is slight and the trajectory is not greatly altered. In the case of two body relaxation the position and velocity (hence energy) of the particles is greatly altered.
\end{quote}

\paragraph{Integrator type.} Typically, the equations of motion are solved using a `Leapfrog' integrator. While this method only gives positions accurate to third order (total error is $ O(h^2)$) it is a symplectic integrator hence it conserves energy (see \ref{sec.sym} for a definition of symplectic). An alternative scheme would be to use the fourth order Runge-Kutta integrator (RK4) which potentially computes positions accurately up to fifth order (but total error is $ O(h^4)$); however, it is not symplectic so does not conserve energy. \citep{bib_nr}

\paragraph{Algorithm speed.} The speed of any algorithm is quoted to ``scale as'' the slowest part (most expensive) of the calculation, in the direct summation case here the slowest part of the algorithm is the force calculation. It is clear from the form of the force calculation that there are $N (N-1)$ force-pairs to be calculated every time-step hence this method is said to ``scale as $N^2$'', where $N$ is the number of particles. In the modern cosmological simulations the above force equation (\ref{eqn_fgr}) for gravity is substituted with the Poisson equation of gravity. This method is outlined in the following subsection. This is designed to provide an improvement on the simulation time.

Direction summation methods have been used to study many astrophysical systems and still prove to be popular. They have been used to study solar system dynamics, planet formation, galaxy formation and galaxy clustering \citep{bib_quin, bib_aar}. Modern Cosmology codes prefer a mix of direction summation (P-P) and the particle-mesh (P-M) method \citep{bib_bert, bib_hmpc}, which are colloquially called $P^3 M$ (``$P$-cubed-$M$").

\subsection{Particle-Mesh $N$-Body}
\label{sec.mesh}
The particle-mesh replaces the force-pair calculations of the direct summation method (above) with a method where the forces upon a particle are calculated with reference to a background gravitational potential. As matter moves it in turn changes the shape of the potential. Particles do not interact directly but via the potential field. They do so using the Poisson equation of gravity, which is equivalent to the Newtonian equation for gravity.

\begin{equation}
\nabla^2 \Phi = 4\pi G a^2 \bar\rho \delta
\end{equation}

\noindent which can be read as the curvature of the gravitational potential $\Phi$ being equivalent to the density $\rho = \bar\rho\delta$ (up to some factor). The above version of the equation is in a typical cosmological form that accounts for expansion via the scale factor $a$. This eventually leads to the calculation of force (or acceleration) that appears in the equations of motion for a cosmological $N$-body code.

The first step in calculating the potential requires the interpolation of particle positions in density. The density field is constructed to be on a regular grid that covers the entire box. The width of the grid spacing (or cell length) is normally that of the interparticle separation in the initial conditions. A typical interpolation scheme is the Triangular Shaped Cloud method (TSC) which appears in the simulation code Hydra \citep{bib_hmpc1, bib_hmpc}. The density is then transformed into Fourier space where it is manipulated to give the transform of potential; crudely $ \Phi_k \sim \rho_k / k^2$.

There are several methods for calculating $\Phi$ in Fourier space, one can crudely divide by $k^2$ or use something more advanced such as a cosine expansion \citep{bib_nr} or via Green's functions, as Hydra does. The cosine expansion is used in the final code that appears in Chapter \ref{ch_5}, some details are given there. Here I will highlight how $k^2$ is replaced by cosines:

\begin{eqnarray}
 \Phi_k &=&  \frac{4\pi G \rho_k}{ 2 \kappa - 3} \nonumber\\
 \kappa &=& (\cos(2 \pi n/L) + \cos(2 \pi m/L) + \cos(2 \pi o/L))
\end{eqnarray}

\noindent $L$ is the number of gridcells per side of the simulation box, $m,n,o$ are indices that run from 0 to $L -1$.

For wave-mechanics it turns out that calculating the potential is sufficient as it appears directly in the Schr\"odinger equation. However, in order to calculate the force a further step is required for $N$-body codes; such codes have to choose an appropriate differencing operator to calculate the force given some potential. Couchman suggests a 10-point differencing operator in the Hydra code, this is more accurate than the usual 2-point operator but far more expensive computationally. A differencing operator is the numerical (or discrete) approximation to the differential operator of calculus, as an example the 2-point forward difference operator is $\Delta u_n = u_{n+1} - u_n$.

\paragraph{Algorithm Speed.} In terms of speed, a P-M code is an improvement over direct force summation: it scales as $O(N \log_2 N)$, so at large $N$ this is considerably faster than $O(N^2)$. The slowest part of the particle-mesh calculation is still the force calculation but now the speed is not due to the pair-wise nature of calculation but due to the inherent speed of the FFT algorithm. Even though modern $N$-body codes rely upon both Particle-Particle interactions and Particle-Mesh interactions, the dominant calculation is still the Particle-Mesh force routine. The use of direct summation is limited to interactions at short distances, while the interactions at longer distances are dealt with using the P-M method described in this section. The use of adaptive space and time stepping can also improve resolution and speed. \citep{bib_aar, bib_nr}

This method forces softening at scales smaller than the gridcell length which truncates the gravitational force. The fundamental limit for force resolution is the Nyquist frequency \citep{bib_hmpc}; a resolution set by the number of gridpoints and hence the total number of wavenumbers in Fourier space. Naturally, this is a drawback for any simulation method that proceeds in this way. The gridcell length is often shortened to just cell length and ``gridpoints'' is often used interchangeably with ``mesh points''.

From this overview of the $P^3 M$ method we believe that some of the key differences and strengths of wave-mechanics should already be apparent. Wave-mechanics requires no interpolation from particle positions to density, so wave-mechanics is always dealing with continuous fields (to machine precision). Secondly, the last step of using a differencing operator to calculate force is not necessary as wave-mechanics uses the potential directly.

The advantage to using $P^3 M$ is that they also use explicit $P$-$P$ calculation (direct summation) when needed. The $P$-$P$ calculation is necessary for calculating force accurately at short distance scales (as the $P$-$M$ method is truncated at short scales). As explained, the $P$-$P$ is slower than it needs to be for larger distance scales but this is the regime when the $P$-$M$ part of a $P^3 M$ is used.

\section{Definition of symplectic}
\label{sec.sym}
The word {\it symplectic} is synonymous with the word {\it Hamiltonian}; therefore, symplectic numerical integrators are specific to Hamiltonian systems and enforce the conservation laws of Hamiltonian dynamics. See \citep{bib_saha} and references therein.

The word symplectic is a construction by Weyl, see his book on Classical groups for the etymology \citep{bib_weyl2}. Now we provide a definition of a symplectic manifold which, in turn, will illustrate how a symplectic integrator works:

\paragraph{Definition} A non-degenerate closed differential 2-form, $\omega$, on $M$ is called a {\it symplectic} (or Hamiltonian) form, and the pair $(M, \omega)$ is a symplectic (or Hamiltonian) manifold. \citep{bib_tens}

To make sense of this we should see how $\omega$ is defined in Hamiltonian dynamics. Using the canonical, or generalized, coordinates for position and momentum ($q, p$; respectively) then $\omega$ (a 2-form) is defined in the following way:

\begin{equation}
\omega =\sum_i \mathrm{d}q^i \wedge \mathrm{d}p_i.
\end{equation}

\noindent here an (exterior) wedge product between the two coordinates is a phase-space volume. Requiring this quantity to be non-degenerate is  merely the formal requirement that there are no zero multipliers: if $\omega = 0$ then either $dq$ or $dp$ must be zero. For real quantities this is always the case (also note that $\omega$ is real by construction).

The definition of closed means that it is divergenceless ($d \omega = 0$). The requirement of $\omega$ to be differentiable means that it is differentiable everywhere, to arbitrary order, on the manifold; hence, the manifold endowed with such a structure is said to be smooth. \citep{bib_tens}

In simple terms this says that the Hamiltonian (energy function) of a system is a smooth real-valued function on a symplectic manifold (commonly referred to as phase space). Given this explanation it becomes quite obvious that any transformation that preserves all of the above can be called a symplectic transformation, or symplectomorphism (which is inherently diffeomorphic). Such a transformation is one of time evolution. Liouville's theorem demonstrates that a symplectomorphism preserves volume in phase-space ($dq \wedge dp$). As the phase-space volume is conserved then the Hamiltonian (total energy of the system) is also conserved.

Of note is that all symplectic manifolds (therefore, groups) are simply connected \citep{bib_weyl2}: hence, all points in phase-space can be continuously transformed from one to another. This is an inherent statement of conservation, the end points are fixed but the path between the two can be continuously deformed however there is only one path: this is much like the variational principle (essentially the principle of least action).

From Noether's theorem we know that the Poincar\'e group, of which the Abelian (symplectic) group is a subgroup, is a fundamental statement about the symmetries and conservation laws of nature. To restate Noether: behind every conservation law is a differentiable symmetry \citep{bib_noe}. It is clear that the topological nature of symplectic manifolds corroborates with the definition of the manifold being smooth and differentiable; it also agrees with the notion of a symplectomorphism. The various mathematical definitions reinforce one another and make sense with what we expect from a physical (Hamiltonian) system. \citep{bib_duff}

Therefore we are inclined to look for a numerical integrator that is symplectic and hence conserves energy. This is an important property for Cosmological codes as the exact positions of particles are less relevant due to the statistical nature of the system. One specific realisation of particle positions is no more important than another, provided they are statistically equivalent (for example, the same degree of clustering: equal values of $\sigma_{8,0}$).

\section{Fluid dynamics and perturbation theory}
\label{sec_pert}
The lack of an (easily computable) analytic solution for the $N$-body codes makes it more difficult to verify whether the simulations are correct. To circumvent this short coming, it is possible to look at the evolution of the Universe using a different technique: perturbation theory. This section looks at the possible advantages of considering a perturbation approach. Similar to the $N$-Body approach, the main idea is to introduce perturbations in density expanded around the mean value. Instead of following a large number of particles this approach evolves the fluid equations (\ref{eqn_fluid}). This can be done in two ways, (1) follow the gravitational collapse of objects within the fluid in configuration space (essentially a mesh based approach), or (2) follow the evolution of the different perturbation modes in Fourier space \citep{bib_mab, bib_huw, bib_cmb, bib_cmb2}.

Most perturbation modes will evolve in a linear manner and hence can be evaluated using Newtonian physics (weak gravitational fields and slow speeds). These equations linearize the fluid equations and hence discard perturbations greater than first order. The procedure for linearizing the fluid equations is found in (for example) Peacock \citep{bib_pea} and in the thesis of Short \citep{bib_short}. Here are the fluid equations in  expanding coordinates:

\begin{eqnarray}
\label{eqn_linflu}
\frac{\partial \delta}{\partial t} + \frac{1}{a} \nabla \cdot v &=& 0 \nonumber \\
\frac{\partial v}{\partial t} + H v &=& - \frac{1}{a} \nabla \Phi \nonumber \\
\nabla^2 \Phi &=& 4\pi G a^2 \bar\rho \delta
\end{eqnarray}

The first line is the continuity equation of the perturbation $\delta$ while the second equation is the Euler equation. The third is the Poisson equation for gravity. The simplest models of perturbation theory often have an analytic solution. One example is the evolution of a CDM-fluid that only includes linear terms in the dynamics: the solution is the well known Einstein de Sitter Universe where $ a \propto t^{2/3}$ and corresponding solutions for $\delta \propto t^n$ are $n = 2/3$ (growing mode) or $-1$ (decaying mode).

The analytic solutions provide a way of calibrating the simulations and also for providing intuition in to harder problems. Given some result from a code that follows non-linear evolution it is necessary to consider what the result means and how to verify it against the real Universe. Such a complex result can be matched with (1) observations and (2) against linear perturbation theory.

A fuller treatment would require studying perturbations using General Relativity (Einstein's field equations). The Boltzmann equation would be recast in a fully relativistic form and all of the constituents of the Universe (radiation, baryonic matter, dark matter, dark energy etc) would combine together to give a full picture. This is difficult computationally but many have tried to write codes that correspond to this fuller picture of perturbation evolution \citep{bib_huw, bib_pea}. Such modifications are needed for following super-horizon perturbations or in regions of strong gravity where a General Relativistic treatment would be necessary. Simpler alternatives have been suggested but they go beyond the scope of this thesis.


\section{From old to new: wave-mechanical approach conceived}
\label{lsswm_con}
The first attempt to describe Cosmic Large Scale Structure using wave-mechanics was by Widrow \& Kaiser \citep{bib_wk} in 1993. They aimed
to overcome the limitations of the previous methods (eg $N$-body, Phase space and Fluid methods). Their goal was to find a model for collisionless
matter that: (1) describes matter as a field rather than particles, (2) is a function only of space and time $(3+1D)$, (3) can follow multiple streams in phase space, and (4) is competitive with the computational time of $N$-body techniques.

Widrow \& Kaiser \citep{bib_wk} provide a summary of each the simulations methods. It is worth re-iterating their comments in order to see why they decided to pursue wave-mechanics (the following comments are paraphrased from their paper).
\newline

$N$-Body (most popular): $N$ `superparticles' with `random' initial positions (determined by the cosmic power spectrum) and velocities. The particles' equation of motion uses simple Newtonian gravity. The positions and velocities approximate the underlying continuous distribution. The aim of cosmological $N$-body codes is to follow the evolution of a large number of particles. Expansion is accounted for by re-writing the equations into comoving coordinates.
\newline

Phase Space method: works directly with distribution function and describes CDM as continuous fluid. It avoids two-body relaxation but has a large number of dimensions and there is difficulty following fine structure in phase space.
\newline

Fluid method: Peebles \citep{bib_peeb2} used a pressureless fluid ($\nabla . v = 0$) in an expanding Universe. He evolved the Euler equation for density and velocity fields (see equation \ref{eqn_linflu}), this ensures mass and momentum conservation. Matter is treated as a continuous field; however, velocity dispersion must be negligible.
\newline

As the $N$-body discrete particles to model a fluid (or something which is similar enough to a fluid) it only approximates the underlying continuous field. Furthermore it requires an additional softening length to prevent two body relaxation. While phase-space methods deal with a continuous distribution function and avoid two-body relaxation they explicitly deal with a single function of a large number of dimensions $(6+1)$, $N$-body methods have proven more tractable when it comes to computer coding. Fluid methods are also continuous; however, the velocity field is single valued, hence cannot handle ``hot'' systems.

The idea proposed by Widrow \& Kaiser is equivalent to the $N$-Body method; however, \citet{bib_sk} developed wave-mechanics in another direction: Schr\"odinger perturbation theory which is an idea similar to the fluid perturbation theory considered in the previous section \ref{sec_pert}. Now that the purpose of why wave-mechanics was chosen is made clear, I will go on to explain what the system of equations mean and how to interpret them in the next chapter.

